\documentclass[10pt,a4paper]{amsart}
\usepackage[margin=19mm]{geometry}
\usepackage{amsmath,amssymb,mathtools}
\usepackage[scr=rsfs,cal=euler]{mathalfa}
\usepackage{xfrac}
\usepackage[unicode,hidelinks,bookmarks=false]{hyperref}
\newcommand{\ie}{i.e.\ }
\newcommand{\cf}{cf.\ }
\newcommand{\resp}{respectively\ }
\newcommand{\Subsection}{\subsection}
\providecommand{\nobreakdash}{\nobreak\hbox{-}}
\setlength{\emergencystretch}{2em}
\multlinegap0pt
\usepackage{tikz}
\usetikzlibrary{cd}
\usepackage{mathtools}
\usepackage{amssymb}
\usepackage[shortlabels]{enumitem}
\usepackage{xcolor}
\usepackage[normalem]{ulem}
\usepackage{float}
\newtheorem{remark}{Remark}
\newtheorem{proposition}{Proposition}
\newtheorem{lemma}{Lemma}
\usepackage{cleveref}
\crefname{remark}{Remark}{Remarks}
\crefname{proposition}{Proposition}{Propositions}
\crefname{lemma}{Lemma}{Lemmas}
\newcommand{\Forest}{\Phi}
\newcommand{\Hom}{\mathrm{Hom}}
\renewcommand{\O}{\mathscr{O}}
\renewcommand{\P}{\mathscr{P}}
\newcommand{\PSh}{\mathrm{PSh}}
\newcommand{\Tree}{\Omega}
\newcommand{\cO}{\O}
\newcommand{\cP}{\P}
\def\colim{\qopname\relax m{colim}}
\newcommand{\el}{\mathrm{el}}
\newcommand{\elem}{\el}
\newcommand{\intarrow}{\rightarrowtail}
\newcommand{\op}{\mathrm{op}}
\title{Day convolution for algebraic patterns}
\author{Thomas Blom, F\'elix Loubaton, Jaco Ruit}
\date{Source: arXiv:2603.29815v2 (CC BY 4.0)}
\begin{document}
\maketitle

\noindent\emph{Source and licence.} Day convolution for algebraic patterns. Version arXiv:2603.29815v2 (CC BY 4.0), distributed by arXiv under the Creative Commons Attribution 4.0 International licence, http://creativecommons.org/licenses/by/4.0/, as recorded in the arXiv licence metadata for this exact version and shown on the abstract page https://arxiv.org/abs/2603.29815v2. Full licence notice: \texttt{licenses/arxiv-2603.29815-CC-BY-4.0.txt}.\par\smallskip

\noindent\textit{Editorial scope.} Counted unit: the article's lemma fib (source0.tex lines 3584--3590). The article's complete original proof of it is delivered with it (source0.tex lines 3591--3637, one \texttt{proof} environment of 47 lines). No mathematics is written, repaired or re-derived: the statement and the proof are the article's own lines, in the article's own notation, and no retained line is substituted or re-flowed.  Retained uncounted as the source-local context the proof needs: lines 678--681; lines 1445--1460; lines 1468--1473; lines 1554--1560; lines 1642--1647; lines 2860--2863; lines 2875--2886.  Nothing was omitted from any retained span, so the package carries no omission record.  No retained line contains a citation, so the wrapper carries no bibliography and no bibliography entry.  The retained lines contain the article's own diagram code, which the wrapper typesets with the standard tikz packages; no image file is delivered and the package declares no asset dependency.  Editorial formatting only, in addition: an AMS \texttt{amsart} preamble that restates the article's own declarations and macros used by the retained lines, namely the environments \texttt{remark}, \texttt{proposition}, \texttt{lemma} on separate counters, together with the standard packages those lines need.  The retained environments print in this document as Remark 1, Proposition 1, Lemma 1 in order of appearance, whereas the article prints the same items with the numbering of its own full document.  The complete native source of the article as distributed in the arXiv e-print for this exact version is bundled unchanged as \texttt{sources/197543/source0.tex}.
\begin{remark}\label{rem:slices-factorization-double-category}
If $\cP$ is a factorization double category, then the Segal condition implies that for every $n \geq 1$, the target map $\cP_n \to \cP_0$ induced by $\{n\} \hookrightarrow [n]$ is a left fibration.
In particular, for any $t = (t_0 \to t_1 \to \cdots \to t_n)$ in $\cP_n$, the target projection $(\cP_n)_{t/} \to (\cP_0)_{t_n/}$ is an equivalence.
\end{remark}

\begin{proposition}\label{prop:RKan-trees-to-forests-simplified}
    Let $\cO$ be an algebraic pattern and $X$ a presheaf on $\Forest[\cO]$ such that for any forest $\langle n;t \rangle$, the canonical map
    \[X(\langle n;t \rangle) \to X(\langle 1;t_{0,1} \rangle) \times_{X(\langle 0;t_1 \rangle)} \cdots \times_{X(\langle 0;t_{n-1} \rangle} X(\langle 1;t_{n-1,n} \rangle)\]
    is an equivalence.
    Then $X$ lies in the image of $i_*$ if and only if for every object $x$ in $\cO$, the map $X(\langle 0;x \rangle) \to \lim_{e \in \cO^\elem_{x/}} X(\langle 0;e \rangle)$ is an equivalence and for any forest $\langle 1;t \rangle$ of length 1, the square
    % https://q.uiver.app/#q=WzAsNCxbMCwwLCJYKFsxO3RdKSJdLFswLDEsIlgoWzA7dF8wXSkiXSxbMSwwLCJ7XFxsaW1cXGxpbWl0c197KHQgXFxpbnRhcnJvdyBzKSBcXGluKFxcY09fMSleXFxlbGVtX3t0L319fSBYKFsxO3NdKSJdLFsxLDEsIntcXGxpbVxcbGltaXRzX3sodCBcXGludGFycm93IHMpIFxcaW4oXFxjT18xKV5cXGVsZW1fe3QvfX19IFgoWzA7c18wXSkiXSxbMCwxLCJkXzEiLDJdLFswLDJdLFsxLDNdLFsyLDMsImRfMSJdXQ==&macro_url=https%3A%2F%2Fuser-uploads.perchance.org%2Ffile%2Fe6181e80ecef834fd03d97c5cca7dae7.bin
    \[\begin{tikzcd}[ampersand replacement=\&]
    	{X(\langle 1;t \rangle)} \& {{\lim\limits_{(t \intarrow s) \in(\cO_1)^\elem_{t/}}} X(\langle 1;s \rangle)} \\
    	{X(\langle 0;t_0 \rangle)} \& {{\lim\limits_{(t \intarrow s) \in(\cO_1)^\elem_{t/}}} X(\langle 0;s_0 \rangle)}
    	\arrow[from=1-1, to=1-2]
    	\arrow["{d_1}"', from=1-1, to=2-1]
    	\arrow["{d_1}", from=1-2, to=2-2]
    	\arrow[from=2-1, to=2-2]
    \end{tikzcd}\]
    is cartesian.
\end{proposition}

\begin{proposition}\label{prop:RKan-trees-to-forest-sound}
    Suppose $\cO$ is sound and let $X$ be a presheaf on $\Forest[\cO]$. Then $X$ lies in the image of $i_*$ if and only if for every forest $\langle n;t \rangle$, the map
    \[X(\langle n;t \rangle) \to \lim\limits_{(t \intarrow s) \in (\cO_{n})^\elem_{t/}} X(\langle n;s \rangle)\]
    is an equivalence.
    If $X$ satisfies the condition from \cref{prop:RKan-trees-to-forests-simplified}, then this only needs to hold for forests of length $\leq 1$.
\end{proposition}

\begin{remark}\label{rem:forest decomposition sound}
    If $\O$ is sound, then it follows from \cref{prop:RKan-trees-to-forest-sound} that the canonical map 
    $$
    \colim_{s\in(\O^{\el,\op}_n)_{/t}} X \boxtimes_{[n]} [n;s] \to X \boxtimes_{[n]} [n;t]
    $$
    is an equivalence for every forest $\left<n;t\right>$ and simplicial space $X$ over $[n]$.
\end{remark}

\begin{proposition}\label{prop:simplicial segal extension}
    Let $\left<n;t\right>$ be a forest. If $X \to Y$ is a (complete) Segal extension between simplicial spaces over $[n]$, then 
    the induced map
    $$X \boxtimes_{[n]} [n;t] \to Y \boxtimes_{[n]} [n;t]$$
    between $\Tree[\O]$-spaces is a (complete) Segal extension as well.
\end{proposition}

\begin{proposition}\label{prop:robust elementary slices decomp}
    Let $\O$ be a robust pattern. If $x\in \O$ lies over $\underline{n}$, then the $\pi_0$-cocartesian inert maps $x\intarrow x^i$ over $\underline{n} \hookleftarrow \{i\}\xrightarrow{=}\{i\}$ induce an equivalence
    $$\coprod_{i=1}^n \O^\el_{x^i/} \to \O^\el_{x/}.$$
\end{proposition}

\begin{remark}\label{rem:decomposition by atoms formula}
Let $\O$ be a robust pattern and  $\left<n;t\right>$ be a forest. Note that the canonical map
    $
    \textstyle \coprod_{i\in \pi_0(t)} (\O_n^{\el})_{t^i/} \to (\O_n^{\el})_{t/}
    $ is an equivalence by \cref{rem:slices-factorization-double-category} and \cref{prop:robust elementary slices decomp},
    where $t\to t^i$ is the cocartesian lift of the map $\pi_0(t_n) \hookleftarrow \{i\} \xrightarrow{=} \{i\}$.
   We may then use the decomposition formula of \cref{rem:forest decomposition sound} to conclude that the canonical map 
    $$
    \coprod_{i\in \pi_0(t)} X \boxtimes_{[n]} [n;t^i] \to X \boxtimes_{[n]} [n;t]
    $$
    is an equivalence in $\PSh(\Tree[\O])$ for every simplicial space $X$ over $[n]$.
\end{remark}

\begin{lemma}\label{lemma:a cartesian fib}
Let $X \to [2;u\star^i t]$ be a complete Segal fibration. Then the induced map
$$
\Hom_{/[2;u\star^i t]}(T_\bullet[1]\boxtimes_{[2]}[2;u\star^i t],X) \to \prod_{i\neq j}\Hom_{/[2;u\star^i t]}([1+\bullet]\boxtimes_{[1]}[1;u_0^j =u_0^j],X)
$$
is map of complete Segal spaces, and can thus be viewed as a functor of categories. As such, it is a cartesian fibration.
\end{lemma}

\begin{proof}
 It follows from \cref{prop:simplicial segal extension} that this is indeed a map between complete Segal spaces.
Note that $[1+\bullet] \simeq T_\bullet[1] \times_{[2]} [1]$, where the pullback is taken along $d_2 \colon [1] \to [2]$.
We get an induced map 
$$\Hom_{/[2;u\star^i t]}(T_\bullet[1]\boxtimes_{[2]}[2;u\star^i t],X)\xrightarrow{\alpha} \Hom_{/[2;u\star^i t]}((T_\bullet[1] \times_{[2]} [1])\boxtimes_{[2]}[2;u\star^i t],X).$$
We claim that this is a right fibration.
To this end, we have to show that the square 
\[\begin{tikzcd}
	\Hom_{/[2;u\star^i t]}({T_1[1]\boxtimes_{[2]}[2;u\star^it]}, X)& \Hom_{/[2;u\star^i t]}({T_0[1]\boxtimes_{[2]}[2;u\star^it]},X)\\
	\Hom_{/[2;u\star^i t]}({(T_1[1]\times_{[2]}[1])\boxtimes_{[2]}[2;u\star^it]},X) & \Hom_{/[2;u\star^i t]}({(T_0[1]\times_{[2]}[1])\boxtimes_{[2]}[2;u\star^it]} , X)
	\arrow[from=1-1, to=1-2]
	\arrow[from=1-1, to=2-1]
	\arrow[from=1-2, to=2-2]
	\arrow[from=2-1, to=2-2]
\end{tikzcd}\]
induced by $\{1\} \to [1]$ is a pullback square.
But this follows from \cref{prop:simplicial segal extension} and the observation that the following square is a pushout in the category of complete Segal spaces over $[2]$:% https://q.uiver.app/#q=WzAsNCxbMCwwLCJUXzBbMV1cXHRpbWVzX3tbMl19WzFdIl0sWzAsMSwiVF8xWzFdXFx0aW1lc197WzJdfVsxXSJdLFsxLDEsIlRfMVsxXSJdLFsxLDAsIlRfMFsxXSJdLFswLDFdLFswLDNdLFszLDJdLFsxLDJdXQ==
\[\begin{tikzcd}
	 {T_0[1]\times_{[2]}[1]} & {T_0[1]} \\
	{T_1[1]\times_{[2]}[1]} & {T_1[1]}
	\arrow[from=1-1, to=1-2]
	\arrow[from=1-1, to=2-1]
	\arrow[from=1-2, to=2-2]
	\arrow[from=2-1, to=2-2]
\end{tikzcd}
\quad \simeq \quad
\begin{tikzcd}
	{[1]} & {[2]} \\
	{[2]} & {[3]}.
	\arrow[from=1-1, to=1-2, "d_2"]
	\arrow[from=1-1, to=2-1, "d_1"']
	\arrow[from=1-2, to=2-2, "d_1"]
	\arrow[from=2-1, to=2-2, "d_3"]
\end{tikzcd}
\]
Now, using the decomposition formula of \cref{rem:decomposition by atoms formula}, the following square is a pullback
% https://q.uiver.app/#q=WzAsNCxbMCwwLCJcXEhvbV97L1syO3VcXGNpcmNeaSB0XX0oWzErXFxidWxsZXRdXFxib3h0aW1lc197WzJdfVsyO3VcXHN0YXJeaSB0XSxYKSJdLFswLDEsIlxccHJvZF97aVxcbmVxIGp9XFxIb21fey9bMjt1XFxjaXJjXmkgdF19KFsxK1xcYnVsbGV0XVxcYm94dGltZXNfe1sxXX1bMTt1XzBeaiA9dV8wXmpdLFgpIl0sWzEsMCwiXFxIb21fey9bMjt1XFxjaXJjXmkgdF19KFsxK1xcYnVsbGV0XVxcYm94dGltZXNfe1sxXX1bMTt0XSxYKSJdLFsxLDEsIjEiXSxbMiwzXSxbMSwzXSxbMCwxLCJcXGJldGEiLDJdLFswLDJdXQ==
\[\begin{tikzcd}
	{\Hom_{/[2;u\star^i t]}([1+\bullet]\boxtimes_{[2]}[2;u\star^i t],X)} & {\Hom_{/[2;u\star^i t]}([1+\bullet]\boxtimes_{[1]}[1;t],X)} \\
	{\prod_{i\neq j}\Hom_{/[2;u\star^i t]}([1+\bullet]\boxtimes_{[1]}[1;u_0^j =u_0^j],X)} & \ast.
	\arrow[from=1-1, to=1-2]
	\arrow["\beta"', from=1-1, to=2-1]
	\arrow[from=1-2, to=2-2]
	\arrow[from=2-1, to=2-2]
\end{tikzcd}\]
As the right vertical morphism is a  cartesian fibration, so is the left vertical one. This implies that $\beta\alpha$ is a cartesian fibration, as desired.
\end{proof}

\end{document}
